\begin{center}
\Huge
	Geometrisk løsning af andengradsligninger
\end{center}
\section*{Bestem rødder geometrisk}
\stepcounter{section}

Det viser sig, at det er muligt at bestemme rødder til andengradspolynomier geometriks Dette kan gøres med passer og lineal på følgende vis. Vi betragter et andengradspolynomium $p$ på formen
\begin{align*}
	p(x) = x^2 +bx + c.
\end{align*}
\begin{enumerate}[label=\roman*)]
	\item Indsæt punktet $A(0,1)$ i et koordinatsystem.
	\item Indsæt punktet $B(-b,c)$ i samme koordinatsystem. 
	\item Indsæt midtpunktet $M$ mellem $A$ og $B$ på koordinatsystemet. Dette har koordinatsættet
	\begin{align*}
		M\left(\frac{-b}{2},\frac{c+1}{2}\right)
	\end{align*}
	\item Tegn en cirkel med centrum i $C$ og som skærer punkterne $A$ og $B$. 
	\item Aflæs cirklens skæringer med $x$-aksen. Disse tilsvarer rødderne for $p$. 
\end{enumerate}

Denne metode til at finde rødder er illustreret på Figur \ref{fig:algoritme}.

\begin{figure}[H]
	\centering
	\begin{tikzpicture}
	\begin{axis}[
	y = 0.5cm, x = 0.5cm,
	axis lines = center, 
	xmin = -4, xmax = 10,
	ymin = -4, ymax = 10,
	xlabel = $x$, ylabel = $y$, 
	ticks = none,
	]
		\draw[color = teal, thick] (axis cs: 2.5,2.5) circle (2.91547*0.5cm);
		\node[color = olive, anchor = east] at (axis cs:0,1) {$A(0,1)$};
		\node[color = olive, anchor = west] at (axis cs:5,4) {$B(-b,c)$};
		\node[color = olive, anchor = south] at (axis cs:2.5,2.5) {$M$};			
		\filldraw[color = olive] (axis cs: 0,1) circle (2pt);
		\filldraw[color = olive] (axis cs: 5,4) circle (2pt);
		\filldraw[color = olive] (axis cs: 2.5,2.5) circle (2pt);
		\node[color = teal, anchor = north] at (axis cs: 1,0) {$x_1$};
		\node[color = teal, anchor = north] at (axis cs: 4,0) {$x_2$};
		\draw[color = teal, thick] (axis cs: 1,0.2) -- (axis cs:1,-0.2);
		\draw[color = teal, thick] (axis cs: 4,0.2) -- (axis cs:4,-0.2);
	\end{axis}
	\end{tikzpicture}
	\caption{Illustration af algoritmen, der bestemmer rødder for andengradspolynomier geometrisk.}
	\label{fig:algoritme}
\end{figure}

Vi ser på et eksempel.
\begin{exa}
	Vi betragter andengradspolynomiet $p$ givet ved
	\begin{align*}
		p(x) = x^2 + 6x + 8.
	\end{align*}
	Vi indsætter punkterne $A(0,1)$ og $B(-6,8)$ i koordinatsystemet og indsætter efterfølgende midtpunktet
	\begin{align*}
		M\left(\frac{-6}{2},\frac{8+1}{2}\right) = M(-3,4.5).
	\end{align*}
	Løsningerne fremgår af Figur \ref{fig:exa} til at være $x = -2$ og $x = -4$.
	\begin{figure}[H]
	\centering
	\begin{tikzpicture}
	\begin{axis}[
	y = 0.5cm, x = 0.5cm,
	axis lines = center, 
	xmin = -10, xmax = 6,
	ymin = -4, ymax = 10,
	xlabel = $x$, ylabel = $y$, 
	]
		\draw[color = teal, thick] (axis cs: -3,4.5) circle (4.60977*0.5cm);
		\node[color = olive, anchor = west] at (axis cs:0,1) {$A(0,1)$};
		\node[color = olive, anchor = east] at (axis cs:-6,8) {$B(-6,8)$};
		\node[color = olive, anchor = south] at (axis cs:-3,4.5) {$M$};			
		\filldraw[color = olive] (axis cs: 0,1) circle (2pt);
		\filldraw[color = olive] (axis cs: -6,8) circle (2pt);
		\filldraw[color = olive] (axis cs: -3,4.5) circle (2pt);
	\end{axis}
	\end{tikzpicture}
	\caption{Eksempel på geometrisk bestemmelse af rødder.}
	\label{fig:exa}
\end{figure}
\end{exa}

\subsection*{Opgave 1}

Bestem rødderne for følgende andengradspolynomier geometrisk. 

\begin{align*}
	&1) \ p(x) = x^2 - 2x - 8 & &2) \ p(x) = x^2 +7x + 10  \\
	&3) \ p(x) = x^2 - 6x + 5  & &4) \ p(x) = x^2 - 3x - 4 \\
	&5) \ p(x) = x^2 + x - 2 & &6) \ p(x) = x^2 + 4x + 3 \\
\end{align*}
